\section{Introduction}
\label{sec:introduction}
Short sensor and keyword sequences provide a useful setting for studying
adaptive Transformer inference on an edge FPGA. Their complete weights can
potentially reside on chip, while projection and feed-forward layers account
for much of the arithmetic. Reducing the tokens that reach later layers can
therefore reduce more than quadratic attention work. Selection, compaction,
summary construction, and mode control also introduce overhead that must be
included in a measured comparison.

\projectname{} restricts adaptation to three finite execution modes in one
compact encoder architecture. It combines W8A8 or W4A8 arithmetic, retention of
chronological groups, a residual summary for each discarded group, and one
mode-selected depth boundary. This scope follows the controlling research
plan's Paper 1 subset; it does not implement a trained energy-budget policy,
head gating, or a general architecture search.

Token removal, merging, and aggregation are established techniques
\cite{arxiv210602034,arxiv220207800,arxiv221009461,arxiv231201026}.
The paper must identify a specific technical distinction and measured benefit
relative to the nearest FPGA work, rather than claim that token summaries or
runtime modes alone are new.
\writingtask{technical lead}{Write the final contribution list after the novelty
review and hardware measurements. State the workload and measurement boundary
for every performance or energy claim.}
